\documentclass[abstract=on,parskip=full,headings=standardclasses,fontsize=11pt,paper=a4]{scrartcl}
%\usepackage[paper=a4paper,left=21mm,right=21mm,top=25mm,bottom=25mm]{geometry}
\usepackage[paper=a4paper,left=22mm,right=22mm,top=20mm,bottom=25mm]{geometry}

\usepackage[utf8]{inputenc}
\usepackage[T1]{fontenc}
\usepackage[english]{babel}

\usepackage{adjustbox}
%\usepackage{amsmath}
\usepackage{graphicx}
%\usepackage{fullpage}
\usepackage{authblk}
\usepackage{setspace}
\usepackage{caption}
\usepackage{booktabs}
\usepackage{url}
\usepackage{comment}
\urlstyle{sf}
\usepackage{lmodern}
\usepackage[parfill]{parskip}
%\usepackage{url}
%\urlstyle{same}
\usepackage[small]{titlesec}
\usepackage{marvosym}

\setcounter{secnumdepth}{0}

\addto\captionsenglish{% Replace "english" with the language you use
  \renewcommand{\contentsname}%
    {Course Structure}%
}



%\deffootnote[10pt]{10pt}{10pt}{\makebox[15pt][l]{\thefootnotemark\hspace{10pt}}}

% Use authoryear-comp to create: (Müller 2015, 2016) instead of (Müller 2015; Müller 2016)

% posscite function

\usepackage[style=authoryear-comp,
    maxcitenames=99,
    maxbibnames=99,
    doi=false,
    %sorting=ynt,
    firstinits=true,
    isbn=false,
    date=short,
    dashed=false,
    url=false,
    sortcites=false,
    backend=bibtex]{biblatex}

\makeatletter
\def\blx@maxline{77}
\makeatother


\DeclareNameFormat{labelname:poss}{% Based on labelname from biblatex.def
  \nameparts{#1}% Not needed if using Biblatex 3.4
  \ifcase\value{uniquename}%
    \usebibmacro{name:family}{\namepartfamily}{\namepartgiven}{\namepartprefix}{\namepartsuffix}%
  \or
    \ifuseprefix
      {\usebibmacro{name:first-last}{\namepartfamily}{\namepartgiveni}{\namepartprefix}{\namepartsuffixi}}
      {\usebibmacro{name:first-last}{\namepartfamily}{\namepartgiveni}{\namepartprefixi}{\namepartsuffixi}}%
  \or
    \usebibmacro{name:first-last}{\namepartfamily}{\namepartgiven}{\namepartprefix}{\namepartsuffix}%
  \fi
  \usebibmacro{name:andothers}%
  \ifnumequal{\value{listcount}}{\value{liststop}}{'s}{}}
\DeclareFieldFormat{shorthand:poss}{%
  \ifnameundef{labelname}{#1's}{#1}}
\DeclareFieldFormat{citetitle:poss}{\mkbibemph{#1}'s}
\DeclareFieldFormat{label:poss}{#1's}
\newrobustcmd*{\posscitealias}{%
  \AtNextCite{%
    \DeclareNameAlias{labelname}{labelname:poss}%
    \DeclareFieldAlias{shorthand}{shorthand:poss}%
    \DeclareFieldAlias{citetitle}{citetitle:poss}%
    \DeclareFieldAlias{label}{label:poss}}}
\newrobustcmd*{\posscite}{%
  \posscitealias%
  \textcite}
\newrobustcmd*{\Posscite}{\bibsentence\posscite}
\newrobustcmd*{\posscites}{%
  \posscitealias%
  \textcites}

\renewbibmacro{in:}{} % no "in" before article

\renewcommand*{\bibpagespunct}{\addcomma\space} % ":" instead of pp
\DeclareFieldFormat{pages}{#1}

% Colon after title
\renewcommand{\subtitlepunct}{\addcolon\addspace }

% Colon instead of pp in references
\renewcommand*{\bibpagespunct}{\addcolon\space}
\DeclareFieldFormat{pages}{#1}

% Colon after name in text
\renewcommand*{\postnotedelim}{\addcolon\space}
\DeclareFieldFormat{postnote}{#1}
\DeclareFieldFormat{multipostnote}{#1}

% Remove brackets around year in bibliography
\usepackage{xpatch,filecontents}

\xpatchbibmacro{date+extrayear}{%
  \printtext[parens]%
}{%
  \setunit*{\addperiod\space}%
  \printtext%
}{}{}

% Supresses URL accessed day
\AtEveryBibitem{%
  \ifentrytype{electronic}
    {}
    {\clearfield{urlyear}\clearfield{urlmonth}\clearfield{urlday}}}
%\DefineBibliographyStrings{english}{%
%urlseen = {Accessed},}

\renewbibmacro*{volume+number+eid}{% number of journal in brackets
 \printfield{volume}%
  %\setunit*{\adddot}% DELETED
  \setunit*{\addnbthinspace}% NEW (optional); there's also \addnbthinspace
  \printfield{number}%
  \setunit{\addcomma\space}%
  \printfield{eid}}
\DeclareFieldFormat[article]{number}{\mkbibparens{#1}}



% Change edition field

\DeclareFieldFormat{edition}%
                   {\ifinteger{#1}%
                    {\mkbibordedition{#1}\addthinspace{}edition}%
                    {#1\isdot}}

% New command to show doi, or url or isbn or issn field
% http://tex.stackexchange.com/questions/48400/biblatex-make-title-hyperlink-to-dois-url-or-isbn
\newbibmacro{string+doiurlisbn}[1]{%
  \iffieldundef{doi}{%
    \iffieldundef{url}{%
      \iffieldundef{isbn}{%
        \iffieldundef{issn}{%
          #1%
        }{%
          \href{http://books.google.com/books?vid=ISSN\thefield{issn}}{#1}%
        }%
      }{%
        \href{http://books.google.com/books?vid=ISBN\thefield{isbn}}{#1}%
      }%
    }{%
      \href{\thefield{url}}{#1}%
    }%
  }{%
    \href{http://dx.doi.org/\thefield{doi}}{#1}%
  }%
}

% Necessary to remove dot after question mark in title
%\newcommand{\killpunct}[1]{}    

% Make full stop after title and before quotation marks in title field
\DeclareFieldFormat{title}{\usebibmacro{string+doiurlisbn}{\mkbibemph{#1}}}
\DeclareFieldFormat[article,incollection,unpublished,phdthesis]{title}%
    {\usebibmacro{string+doiurlisbn}{\mkbibquote{#1}}}
   % {\usebibmacro{string+doiurlisbn}{\mkbibquote{#1.\isdot}}}

\renewcommand*{\newunitpunct}{.\space}


\bibliography{/Users/smueller/Documents/GitHub/literature/muellerlibrary.bib}
%\bibliography{/Users/stefan/GitHub/literature/muellerlibrary.bib}


\usepackage{xcolor}
%\definecolor{JournalBlue}{RGB}{0, 12, 146}

\definecolor{JournalBlue}{RGB}{25, 63, 144}

%https://en.wikibooks.org/wiki/LaTeX/Colors
\usepackage[colorlinks=true, linkcolor=JournalBlue, filecolor=black, urlcolor=JournalBlue, pdfborder={0 0 0},citecolor=JournalBlue]{hyperref}%RoyalBlue
%\usepackage[colorlinks]{hyperref}

\clubpenalty = 10000 
\widowpenalty = 10000 
\displaywidowpenalty = 10000

\setlength\parindent{0pt}


\usepackage{titlesec}
\titleformat{\section}
{\normalfont\large\bfseries\onehalfspacing}{\thesection}{3em}{}

\titleformat{\subsection}
{\normalfont\bfseries\onehalfspacing}{\thesection}{3em}{}




\begin{document}
	

\singlespacing

\begin{comment}
\noindent
\adjustbox{valign=t}{\begin{minipage}{0.24\textwidth}% adapt widths of minipages to your needs
%\includegraphics[width=\linewidth]{pictures/ucd_logo}
\includegraphics[width=\linewidth]{pictures/ucd_cpl_merged.png}
\end{minipage}}%
\hfill%
\adjustbox{valign=t}{\begin{minipage}{0.9\textwidth}\raggedleft
{%\footnotesize
\textbf{Stefan Müller, PhD} \\
Assistant Professor and Ad Astra Fellow \\
School of Politics and International Relations\\
University College Dublin \\
Belfield, Dublin 4, Ireland \\
%\Telefon\ + 353\,89\,975\,25\,79 \\
\Letter\ \href{mailto:stefan.mueller@ucd.ie}{stefan.mueller@ucd.ie}\\
\ComputerMouse\ \url{https://muellerstefan.net} \\
}
\end{minipage}}
\end{comment}


\noindent
\adjustbox{valign=t}{\begin{minipage}{0.3\textwidth}% adapt widths of minipages to your needs
\includegraphics[width=0.8\textwidth]{pictures/ucd_cpl_merged}
\end{minipage}}%
\hfill%
\adjustbox{valign=t}{\begin{minipage}{0.7\textwidth}\raggedleft
{%\small 
\textbf{Stefan Müller, PhD}\\
Assistant Professor and Ad Astra Fellow \\
School of Politics and International Relations\\
University College Dublin \\
Belfield, Dublin 4, Ireland \\
\Letter\ \href{mailto:stefan.mueller@ucd.ie}{\textsf{stefan.mueller@ucd.ie}}\\
\ComputerMouse\ \url{https://muellerstefan.net} \\
}
\end{minipage}}



\singlespacing
\vspace{1cm}

\begin{center}
{\large %Year 3 Module; 
Spring Trimester 2021} \\ 
\bigskip

{\Large \textbf{Connected\_Politics} (\href{https://sisweb.ucd.ie/usis/!W_HU_MENU.P_PUBLISH?p_tag=MODULE&MODULE=POL42350}{POL42350})} 
\bigskip


{\large  {Last update: \today}}\\
\bigskip

%\textbf{Note: Draft -- subject to change}

Latest version: \url{https://muellerstefan.net/teaching/2021-spring-connected-politics.pdf}
\end{center}

\vspace{1.5cm}



\hrule
\medskip
% first column
\begin{minipage}[t]{0.5\textwidth}
Term: Spring Trimester 2021 \\
Time: Wednesday, 2:00pm--3:50pm \\
Location:  online (Zoom) \\ %\href{https://www.ucd.ie/t4cms/UCD%20Student%20Centre.pdf}{Lecture: QUI 006; Lab: F20 Newstead} \\
ECTS: 10.0 \\
Format: group work; seminar attendance
\end{minipage}
%second column
\begin{minipage}[t]{0.49\textwidth}
\begin{flushright}
Module Coordinator: Stefan Müller \\
 \href{mailto:stefan.mueller@ucd.ie}{\textsf{stefan.mueller@ucd.ie}} \\
 \url{https://muellerstefan.net} \\
Office:  Newman Building \\
Office hours: Tuesday, 10:00--11:30 (via Zoom)
\end{flushright}
\end{minipage}
\medskip
%\vspace{2.5mm}
\hrule 

\section*{Introduction}


Welcome to Connected\_Politics! This module trains you to conduct research projects relating to computational social science in small teams under the supervision of an assigned project coordinator and the module coordinator. You will apply cutting-edge methods, such as quantitative text analysis, machine learning, image recognition, and network analysis, to answer social science research questions. You will learn how to collaborate on research projects with your peers, setting out short-term and longer-term goals, and dividing up various tasks within  groups. At the end of the module, you will have gained significant experience in designing and executing a collaborative academic research project.



\begin{comment}
For the applied parts of this course, such as data import, data wrangling, and data visualisation, we will read parts of the following textbooks. Both books will also help you with your homework assignments and the technical elements of your course paper. 

\begin{itemize}
\item \fullcite{wickham17}. URL: \url{https://r4ds.had.co.nz}.
\item \fullcite{healy19}. URL: \url{https://r4ds.had.co.nz}.
\end{itemize}


We will work extensively with the \textsf{R} statistical programming language. The three books mentioned above \autocite{ismay20,wickham17,healy19}   provide detailed and intuitive examples and the corresponding \textsf{R} code (based on the \href{https://www.tidyverse.org}{\texttt{tidyverse}} approach). In addition to these books, I recommend the following literature for introductions to statistical methods, regression, causal inference, and \textsf{R}:

\begin{itemize}
\item \textbf{Basic grasp of statistics}: \fullcite{spiegelhalter20}.
\item \textbf{Research design}: 
\begin{itemize}
\item \fullcite{kellstedt18}.
\item \fullcite{king94}.
\end{itemize}
\item \textbf{\textsf{R} and regression analysis}: 
\begin{itemize}
\item \fullcite{larsen}.
\item \fullcite{bryan}.
\item \fullcite{gelman20}.
\end{itemize}
\item \textbf{Data visualisation}: \fullcite{wilke19}.
\item \textbf{Causal inference}: 
\begin{itemize}
%\item \fullcite{buenodemesquita}.
\item \fullcite{cunningham}.
\end{itemize}
\item \textbf{RMarkdown}: \fullcite{wickham17}.
\end{itemize}

\end{comment}


\section*{Learning Outcomes}

\begin{enumerate}
\item  Execute a demanding research project using methods relating to computational social science
\item  Collaborate with peers and academic faculty on an academic research project
\item Evaluate and compare a variety of research methods, sources, data, and analysis
\item Critically and thoroughly examine a research question through independent, data-driven research
\item Effectively communicate methods and findings
\end{enumerate}


\section*{Indicative Module Content}


\begin{itemize}
\item Working on collaborative projects
\item Research design(s) and the role of theory in the ``digital age''
\item Formulating and designing a research question
\item Case-selection strategies
\item Operationalisation and measurement
\item  Open science practices, research transparency in groups
\item Replicability and reproducibility of research
\item Presentation of progress
\end{itemize}


\section*{Approaches to Teaching and Learning}


This project will train you how to comment critically and constructively on working papers during and after  research seminars, and how to conduct a demanding research project using methods relating to computational social science. To reach these goals, you will  attend the Connected\_Politics Lab seminar series,  work in groups,  allocate tasks, present your progress, and write a research paper. The module centres on active and task-based learning in groups along with interactive seminar discussions. In addition, we will have check-in meetings (around 30 minutes) after each seminar to discuss the progress and open questions.




\section*{Expectations, Assessment, and Grading}


\begin{itemize}
\item \textbf{Seminar} (throughout the trimester): Attending the \href{https://www.ucd.ie/connected_politics/events/}{Connected\_Politics Lab} seminar in Spring Term 2021 and writing response papers on presentations (pass/fail; at least 5 response papers must receive a `pass'  grade); response papers must be submitted until the Monday (8:00pm)\footnote{All times in this syllabus are listed in \href{https://www.timeanddate.com/time/zones/ist-ireland}{Irish Standard Time (IST)}.} after  the presentation [20\% of final grade]
\item \textbf{Presentation} (week 7):  Conference-style presentation of the research question, data, methods, initial results, and progress on the project  on Wednesday, 2 March 2021 [20\% of final grade]
\item \textbf{Group Project} (week 12): A 6,000-word research paper; deadline Tuesday, 4 May 2021, 8:00pm, IST	  [60\% of final grade]
\end{itemize}




%\section*{Expectations and Grading}

\subsection*{Response Papers}


Students attend the \href{https://www.ucd.ie/connected_politics/events/}{Connected\_Politics Lab Seminar Series} and write \textbf{one-page response papers} on the presentations by external presenters. First, you should summarise the research project in 2--3 sentences. Afterwards,  you should identify either a limitation of the project or a possible extension. Note that what you propose should be feasible (ideally by you). If, for example, you find the author's data weak, then you should identify better data, or at least propose a plausible way of collecting these data. If you think the method is wrong, explain why and suggest a better one. If the conclusions do not follow from the premises, discuss what conclusions are actually supported.  You should outline a specific course of action.  
 
 
Students need to attend the seminars listed below, and submit at least five response papers with a ``pass grade''. The Zoom links will be distributed on Brightspace and the Slack workspace for this module.  Response papers must be submitted until the Monday (8:00pm) after the presentation. For example, for the presentation on 27 January (Wednesday), the response paper must be submitted no later than Monday, 1 February, 8:00pm.


All seminars take place via Zoom between 2pm and 2:50pm. After each presentation, we will have an informal check-in meeting to discuss questions that came up during your group work.

\begin{itemize}
\item 27 January 2021 -- Jesper Lindqvist (UCD): \textit{A Political Esperanto, or False Friends? – 'Left' and 'Right' in Different Political Contexts}
\item 10 February 2021 -- Kevin Munger (Penn State University): \textit{Fifteen Seconds of Fame: TikTok and the Democratization of Mobile Video on Social Media}
\item 24 February 2021 -- Olessia Koltsova (HSE University): \textit{What Do Online Experiments Tell Us About Political Fake News Recognition and Trust?}
\item 24 March 2021 -- Sandra González-Bailón (University of Pennsylvania): \textit{Exposure to News in the Digital Age: How Online Networks Shape the Consumption of Political Information}
\item 7 April 2021 -- Anita Gohdes (Hertie School of Governance): \textit{Online and Offline Responses to Protest in Electoral Autocracies}
\item 21 April 2021 -- Taha Yasseri (UCD): \textit{The Double-edged Sword of Online Politics}
\end{itemize}


\subsection*{Group Work}

In the introductory session on 20 January 2021 (2:00pm--3:50pm), project coordinators will present a variety of research projects. You can choose three projects and you will be allocated to one of these projects. You work with your group on this project throughout the entire term. 


\subsubsection*{Group Presentations}
\label{presentations}
Each group \textbf{presents the progress} of their research projects in Week 7 (Wednesday, 3 March, 2:00pm--4:00pm, IST).% This presentation should outline the research question, the theoretical expectations for your empirical analysis, describe the data, provide descriptive statistics and plots of the data, and (if possible) include a preliminary analysis. 


Each presentation should last between 10 and 12 minutes. Make sure that your presentation does not exceed 12 minutes (we may stop your presentation after 12 minutes!), and practice the presentation with your group in advance.  Please prepare slides (LaTeX or PowerPoint/Google Docs) and send me the slides no later than 10am on Wednesday, 3 March. Note that we will give the same grade to all group members.

The presentation should present your progress by covering the following aspects:

\begin{enumerate}
\item What is your research question and why is it important?
\item What are your theoretical expectations? Which main hypothesis are you going to test in your research paper?
\item Which data are you going to use? Provide descriptive statistics and a short overview (e.g. number of documents/texts/tweets; time span etc.). Describe the data in one or two graphs that effectively communicate the nature of the data source.
\item Which method are you going to apply to test your hypothesis?
\item Ideally, provide initial results.
\item What issues did you have encountered or what are the risks associated with the next steps (if applicable)? 
\end{enumerate}

We encourage presentations to be shared between two and three students (depending on the group size).  
After each presentation, we will have a Q\&A which will last for around 10--15 minutes. You will receive questions from your peers and the project coordinators. Students who do not present are responsible for providing answers to the questions. Therefore, we strongly encourage you to divide up tasks, practice the presentation,  and discuss questions that may come up.


Kieran Healy has published an excellent blog post\footnote{\url{https://kieranhealy.org/blog/archives/2018/03/24/making-slides/}.} on how to give  presentations and how to design slides that effectively communicate your work and results. We encourage you to follow his advice.

The grade will be based on your progress, the content and clarity of the presentation, and  your ability to reply to the questions. 

\subsubsection*{Research Paper}

 Each group submits a \textbf{6,000 word research paper}. The research paper builds on the proposals by the project coordinators and the presentation in Week 7. You  must submit the research papers no later than \textbf{Tuesday, 4 May 2021} (8:00pm). Each group will receive one grade (not individual grades per student) for the presentation and research paper, but you will be asked to clearly indicate who took over which parts of the project. 
 
 

The research paper should contain the following sections:
\begin{itemize}
\item \textbf{Introduction and research question} 
\begin{itemize}
\item Explain the puzzle and research question
\item Highlight the relevance
\item Include the central hypothesis to be tested
\end{itemize}
\item \textbf{Theory and expectations} 
\begin{itemize}
\item Explain the theoretical assumptions based on previous findings regarding the  relationship between your dependent variable and the key independent variable
\end{itemize}
\item \textbf{Methodology} 
\begin{itemize}
\item Describe your dataset, the unit of analysis, the number of observations included in the analysis, the number of missing observations (if appropriate), the measurement of key variables, and the empirical analysis
\end{itemize}
\item \textbf{Results}
\begin{itemize}
\item Present the results of your empirical analysis.
\end{itemize}
\item \textbf{Conclusion}
\begin{itemize}
\item Referring back to the introduction, what can we conclude, and what have we learned?
\end{itemize}
\end{itemize}



Feedback will be provided by your project coordinator and the module coordinator  throughout the module. The module coordinator (Stefan Müller) will grade the response papers, presentations, and research papers. All assignments will be uploaded on Brightspace.  The communication for this module takes place through Slack. Please make sure to check the workspace at least once a day. I also recommend downloading the Slack desktop app.\footnote{\url{https://slack.com/intl/en-ie/downloads}.}


Each group will receive one grade (not individual grades per student) for the presentation and the research paper. However, you will be asked who took over which parts of the project.  More details on the presentation will be provided in a separate document. 

\textit{Important}: it is the students' responsibility to raise alarm if collaboration in your project work is lacking. Please inform the module coordinator as soon as possible. Moreover, it is your responsibility to schedule regular group meetings (1--2 per week).

The presentations and research paper will be graded according to \href{https://bit.ly/3bLcgRp}{UCD's Module Grade Descriptors}.\footnote{\url{https://bit.ly/3bLcgRp} (PDF).}

\begin{comment}

For information on academic writing, I recommend the following two sources:
\begin{itemize}
\item \fullcite{heard16}.
\item \fullcite{dunleavy14}.
\end{itemize}

If you require information on proper citation style, please refer to the guidelines of the American Political Science Association:

\begin{itemize}
\item \fullcite{apsa18}.
\end{itemize}
\end{comment}



\begin{table}[h] \centering \onehalfspacing
\caption*{Student effort hours}
\begin{tabular}{ l r} 
\toprule
Student effort type &  Hours \\
\midrule
Lectures & 13 \\
Autonomous Student Learning  & 211 \\
\textbf{Total} & \textbf{224} \\
\bottomrule
\end{tabular}
\end{table}


\section*{Plagiarism}

Although this should be obvious, plagiarism -- copying someone else's text without acknowledgement or beyond `fair use' quantities -- is not allowed. Plagiarism is an issue we take very seriously here in UCD. Please familiarise yourself with the definition of plagiarism on UCD's website\footnote{\url{https://libguides.ucd.ie/academicintegrity}.} and make sure not to engage in it.



\section*{Late Submission Policy}

All written work must be submitted on or before the due dates. Students/groups will lose one point of a grade for work up to 5 working days late ($B-$ becomes $C+$). Students will lose two grade points for work between 5 and 10 working days late ($B-$ becomes $C$). When more than two weeks are necessary, the student will need to apply for extenuating circumstances application via the SPIRe Programme Office.




\section*{Questions and Problems}

In this module, we will discuss concepts, methods, and software you might not have heard of before. I am  aware that parts of this module could be challenging and we will assist you as best as I can. In addition to the meetings after Connected\_Politics seminars, I  offer weekly office hours for  participants of this module. My office hours will take place via Zoom on Tuesday between 10am and 11:30am. I will share the link and password to the virtual room in the first lecture and post it on Brightspace. 

If you struggle to solve problems relating to \textsf{R}, Python, specific software packages or statistical concepts, please follow the steps outlined below prior to contacting your project coordinator and/or the module coordinator. %It is very likely that at least one other person faced the same problem before or received the same error message. % If something is not working, try the following steps first before contacting me: 
\begin{itemize}
\item Check the literature provided by your project coordinator and the literature suggestions on Brightspace.
\item Get in touch with the other members of your group and try to solve the issue together.
\item For questions about software packages: read the package documentation and/or the documentation of a  function that returns the error.
\item For questions about statistical concepts, I recommend to consult the \href{https://hbiostat.org/doc/glossary.pdf}{Glossary of Statistical Terms}, \textcite{vanatteveldt21}, \textcite{ismay20}, \textcite{wickham17}, and \textcite{healy19}.
\item Try to summarise the problem in your own words and then google this summary. If the problem relates to \textsf{R}, add \texttt{rstats} to your search query; for questions about Python add \texttt{python}. For example: \texttt{how to import csv file in rstats}. I am almost certain that you find a solution to most of your questions. 
\item If your code returns an error, google the text of the error message.  For example: \texttt{Error: Can't subset columns that don't exist}.
\end{itemize}

$\longrightarrow$ If the steps above still do not solve your problem or question, please ask your question in the \texttt{\#questions}  channel on Slack. We are happy to help, but keep in mind that coding and collaborative work always involves trial and error, and that `learning by doing' is essential and an integral part of this module.

\section*{Syllabus Modification Rights}

I reserve the right to reasonably alter the elements of the syllabus at any time by adjusting the reading list to keep pace with the course schedule. Moreover, I may change the content of specific sessions, depending on the participants' prior knowledge and research interests. If I make adjustments, I will send an email to all seminar participants and upload the revised syllabus to Brightspace.

%\newpage



\tableofcontents

\section{20 January 2021: Introductory Session} \label{introsession}

In our introductory session, we will outline the content and expectations for this module. Besides, the project coordinators will shortly present  research projects for this module. After this session, we will distribute an online form where students can express up to four preferences. The following texts (all on Brightspace) offer a good and concise introduction to the field of computational social science. Please read these texts before class:


\begin{itemize}
\item \fullcite{lazer20}. 
\item \fullcite[ch.1--2]{salganik18} (focus mainly on ch. 1; skim ch. 2).
\item \fullcite[ch. 1]{vanatteveldt21}.
\end{itemize}
%\newpage
%\sloppy
%\renewcommand*{\bibfont}{\small}
%\onehalfspacing
%\setlength{\bibitemsep}{0.2em} % increase space between references
%\printbibliography



\section{27 January 2021 -- Jesper Lindqvist (UCD)}

\textit{A Political Esperanto, or False Friends? – `Left' and `Right' in Different Political Contexts} (with Jos Elkink)
 
\textit{Abstract}: The Left-Right dimension has been and continues to be a prominent component of advanced representative democracies, which is used to simplify the political landscape. Nevertheless, it is unclear why the same two terms are used in multiple countries. This would imply that the terminology has a similar core meaning in different political contexts. Yet no such stable element has been established in the political science literature. This paper examines five different possible criteria that have been proposed to separate left from right: change/resistance to change, secular/religious, equality/inequality, equality of outcome/equality of opportunity and for/against government intervention in the economy. We examine these criteria in eight different countries (with varying political contexts), by studying responses to open-ended survey questions on what the terms ``left'' and ``right'' mean. The data are analysed using quantitative text analysis (more specifically topic modelling through Non-negative Matrix Factorization) to examine how respondents understand the left-right terminology. The overall results demonstrate varied support for the different explanations, with the most support found for equality/inequality, for/against government intervention in the economy, as well as change/resistance to change. We find little evidence for the two criteria secular/religious and equality of outcome/equality of opportunity.


\textbf{About the speaker}: Jesper Lindqvist is a PhD Candidate in the School of Politics and International Relations at University College Dublin. His current research is focused on understanding the meaning of left-right politics in representative democracies. In addition, his research interests also include public opinion, ideological dimensions and democratic representation.


We will have an informal check-in meeting after the seminar (from 3:00pm--3:30pm) to discuss questions that came up during your group work.

\section{10 February 2021 -- Kevin Munger (Penn State University)}
 
 
\textit{Fifteen Seconds of Fame: TikTok and the Democratization of Mobile Video on Social Media}


\textit{Abstract}: TikTok has rapidly developed from a punchline for jokes about “kids these days" into a formidable force in American politics. The speed of this development is unprecedented, even in the rapidly-changing world of digital politics. Through a combination of hashtag and snowball sampling, we identify 5,495 TikTok accounts who primarily post about politics, allowing us to analyze trends in the posting, viewing and commenting behavior on 712,193 tiktoks they have uploaded. We test a number of theories about how the unique combination of a ordances on TikTok shapes how it is used for political communication.

\textit{About the speaker}: Kevin Munger is an Assistant Professor of Political Science and Social Data Analytics at Penn State University. His research looks at social media and other contemporary internet technology has changed political communication. Kevin has published research on the subject using a variety of methodologies, including textual analysis, field experiments, longitudinal surveys and qualitative theory. Kevin's research has appeared in leading journals like the American Journal of Political Science, Political Behavior, Political Communication, and Political Science Research and Methods. His present interests include cohort conflict in American politics and developing new methods for social science in a rapidly changing world.
  
  
  
We will have an informal check-in meeting after the presentation (from 3:00pm--3:30pm) to discuss questions that came up during your group work.
  
  
\section{24 February 2021 -- Olessia Koltsova (HSE University)} 

\textit{What Do Online Experiments Tell Us About Political Fake News Recognition and Trust?}



\textit{About the speaker}: Olessia Koltsova is a social scienist committed to interdisciplinary research in the sphere of computational communication science. As the director of Laboratory for Social and Cognitive Informatics (formerly – Laboratory for Internet Studies), she leads various collective projects in the sphere of internet and society, as well as in methods of large-scale automatic internet data analysis for social science. In recent years, she has published on online community structure, cognitive and psychological limits to online communication, ethnicity-targeted speech detection, user content topical composition and sentiment, and other topics. She is also the author of News Media and Power in Russia, Routledge, 2006.
  
  
We will have an informal check-in meeting after the presentation (from 3:00pm--3:30pm) to discuss questions that came up during your group work.


\section{3 March 2021 -- Group Presentations}

In this session, each group will present their progress. Detailed instructions about the presentation and format are outlined on page \pageref{presentations}.

%More details on the structure of the presentations will be provided on Brightspace.

\section{24 March 2021 -- Sandra González-Bailón (University of Pennsylvania)}


 \textit{Exposure to News in the Digital Age: How Online Networks Shape the Consumption of Political Information}
 

\textit{Abstract}: The abundance of media options is a central feature of today’s information environment, and information simultaneously flows through many parallel channels (i.e., the web, social media, TV). In this talk I will discuss recent research in which we measure exposure to news across channels to (1) test claims of increasing audience fragmentation and ideological segregation and (2) measure the influence of automated accounts in distorting the salience of news sources on social media. Using an unprecedented combination of observed data from the US comprising a five-year time window and involving tens of thousands of panelists, I will show that co-exposure to diverse news is on the rise. And using social media data from two contentious political events in France and Spain, I will show that verified accounts are significantly more visible than unverified bots, and that discrepancies in source salience in social media and the web are generated by both human and bot activity. I will discuss the implications of these findings for how we think about the current communication environment, exposure to news, and ongoing attempts to limit the effects of misinformation, including social media verification policies.

\textit{About the speaker}: Sandra González-Bailón is an Associate Professor at the Annenberg School for Communication at the University of Pennsylvania, and affiliated faculty at the Warren Center for Network and Data Sciences(link is external). Prior to joining Penn, she was a Research Fellow at the Oxford Internet Institute(link is external) (2008-2013). She completed her doctoral degree in Nuffield College(link is external) (University of Oxford) and her undergraduate studies at the University of Barcelona(link is external). Her research lies at the intersection of network science, data mining, computational tools, and political communication. Her applied research looks at how online networks shape exposure to information, with implications for how we think about political engagement, mobilization dynamics, information diffusion, and news consumption. 
  
  
 
 We will have an informal check-in meeting after the presentation (from 3:00pm--3:30pm) to discuss questions that came up during your group work.
 
\section{7 April 2021 -- Anita Gohdes (Hertie School of Governance)}

 \textit{Online and Offline Responses to Protest in Electoral Autocracies}


\textit{Abstract}: An influential scholarship shows that contemporary autocrats respond to dissent using a combination of online and offline tactics. Yet, we know little about the relationship between online and offline responses to protest, in particular in the context of electoral autocracies. In this paper we study the extent to which responses in the online and offline sphere operate as complements or substitutes. In addition, we explore whether the regime’s uncertainty with regard to the scale of protests affect the synchronization of online and offline tactics. To gain traction on these question, we investigate the Russian government’s response to two waves of large-scale, anti-corruption protests taking place in March and June 2017. Offline, we analyze strategies of preventive repression, such as officially denying permissions for demonstrations, as well as responsive repression, such as arrests of protesters. Online, we study the social media activity of accounts belonging to the Prime Minister and the Kremlin. Our preliminary descriptive results show important variation in online and offline strategies adopted in anticipation of, and in response to the March and June protests. They also suggest that in both instances, online responses followed similar patterns to the offline repression. Our findings have important implications for our understanding of authoritarian resilience in the context of increasingly networked protest.


\textit{About the speaker}: Anita Gohdes is a Professor of International and Cyber Security at the Hertie School in Berlin. Previously, I was Assistant Professor of International Relations at the University of Zurich, and postdoctoral research fellow at the Belfer Center and the Women and Public Policy Program in the Harvard Kennedy School. Her research focuses on political violence, state repression and the measurement of human rights. Anita's work has been covered by various news outlets and is accepted or appears in the American Journal of Political Science,  Journal of Politics, Journal of Conflict Resolution, Journal of Peace Research, Significance, Journal of Human Rights, and at Oxford University Press. 

  
  

We will have an informal check-in meeting after the presentation (from 3:00pm--3:30pm) to discuss questions that came up during your group work.


\section{21 April 2021 -- Taha Yasseri (UCD)}


 \textit{The Double-edged Sword of Online Politics}


\textit{Abstract}: An influential scholarship shows that contemporary autocrats respond to dissent using a As people go about their daily lives using social media, such as Twitter and Facebook, they are invited to support myriad political causes by sharing, liking, endorsing, viewing and following. Chain reactions caused by these tiny acts of participation form a growing part of collective action today, from neighbourhood campaigns to global political movements. On the other hand these very activities produce large scale transactional data that are unprecedented in political science studies and provide the researchers with unique opportunities to study such collective actions.This lecture shows how most attempts at collective action online fail. Those that succeed can do so dramatically, but are unpredictable, unstable, and often unsustainable. The presented research is based on application of social network analysis, data and text mining and data science techniques to large scale transactional data generated on social media and other internet-based platforms.


\textit{About the speaker}: Taha Yasseri is an Associate Professor at the School of Sociology and a Geary Fellow at the Geary Institute for Public Policy at University College Dublin, Ireland. Formerly he was a Senior Research Fellow in Computational Social Science at the Oxford Internet Institute, a Turing Fellow at the Alan Turing Institute for Data Science, and a Research Fellow in Humanities and Social Sciences at Wolfson College, University of Oxford. Taha Yasseri has interests in analysis of large-scale transactional data and conducting experiments to understand human dynamics, government-society interactions, mass collaboration and collective intelligence, information and opinion dynamics, collective behaviour, and online dating.


We will have an informal check-in meeting after the presentation (from 3:00pm--3:30pm) to discuss questions that came up during your group work.




\section{Additional Meetings Throughout the Term}

Besides the \href{https://www.ucd.ie/connected_politics/events/}{six Connected\_Politics Lab seminars}, the initial meeting on 20 January, and the group presentations on 3 March, you will be meet the other group members (virtually) to work on  the research project. It is your task to organise group meetings, allocate tasks, and communicate with your peers. We recommend that you communicate and chat through \href{https://ucdconnectedpolitics.slack.com.}{Slack} and that you have at least one or two  group meetings  per week to discuss your progress and allocate  tasks. 


 In addition, you will have at least two meetings with your project coordinator. The project coordinators will provide a list of literature to get you started with your project, answer selected questions about methods or software. Yet, it is your task to get familiar with prior work, the required software packages and writing code -- working with software and coding collaboratively are integral parts of the group work. The project coordinator will also meet you after your presentation to discuss strengths and weaknesses and the required actions for the research paper.


\textit{Important}: it is the students' responsibility to raise alarm if collaboration in your project work is lacking. Please inform the module coordinator as soon as possible.


If anything is unclear, you can always contact the module coordinator (Stefan Müller) on  Slack.


\end{document}








